\section*{Capítulo \ref{ch:b3:asymmetric-synthesis} --- Síntesis asimétrica}

\begin{solution}{exo:b3:asymmetric-synthesis:1}
ee del 80\,\%: $\mathrm{er} = 1.8/0.2 = 9$ (90\,:\,10). er de 99\,:\,1: $\mathrm{ee} = 98/100 = 98$\,\%. ee del 50\,\%: 75\,\% y 25\,\%.
\end{solution}

\begin{solution}{exo:b3:asymmetric-synthesis:2}
$(1840 - 60)/(1840 + 60) = 0.937$: ee del 93.7\,\%, una er de $1840/60 = 31$.
\end{solution}

\begin{solution}{exo:b3:asymmetric-synthesis:3}
Prioridades: O $>$ etilo $>$ H. Dibujado con el O arriba, el etilo abajo a la izquierda y el H abajo a la derecha, la
cara delantera es \textit{Si}, y la trasera, \textit{Re}. Un grupo metilo añadido por la \omterm{def:b3:asymmetric-synthesis:faces}{cara Re}, con
las prioridades OH $>$ etilo $>$ metilo $>$ H en el producto, da el (\textit{R})-butan-2-ol.
\end{solution}

\begin{solution}{exo:b3:asymmetric-synthesis:4}
Etanol: \omterm{def:b3:asymmetric-synthesis:topicity}{enantiotópicos} (los intercambia el \omterm{def:b3:point-groups:mirror}{plano de simetría} de la molécula; mismo
desplazamiento en un disolvente aquiral). Butan-2-ol: \omterm{def:b3:asymmetric-synthesis:topicity}{diastereotópicos} (contiguos a un estereocentro;
desplazamientos distintos).
\end{solution}

\begin{solution}{exo:b3:asymmetric-synthesis:5}
er de 199: $RT\ln 199 = \qty{13.1}{kJ/mol}$ a \qty{298}{K}, y \qty{8.6}{kJ/mol} a \qty{195}{K}. Un ee del 80\,\% a \qty{25}{\celsius} es una er de 9,
$\Delta\Delta G^\ddagger = RT\ln 9 = \qty{5.45}{kJ/mol}$; a \qty{-78}{\celsius}, $\mathrm{er} = \eu^{5450/(8.314 \times 195.15)} = 28.7$, ee del 93\,\%.
\end{solution}

\begin{solution}{exo:b3:asymmetric-synthesis:6}
\omterm{def:b3:asymmetric-synthesis:ee}{Pureza óptica}: $12.0/15.0 = 80$\,\%: 90\,\% de (\textit{S}) y 10\,\% de (\textit{R}). Una rotación pequeña, las impurezas,
los efectos de concentración y de disolvente y una relación rotación--composición no lineal
la hacen poco fiable; la cromatografía separa y cuenta directamente los
enantiómeros.
\end{solution}

\begin{solution}{exo:b3:asymmetric-synthesis:7}
$s = \ln(0.40 \times 0.10)/\ln(0.40 \times 1.90) = -3.22/-0.274 = 12$.
\end{solution}

\begin{solution}{exo:b3:asymmetric-synthesis:8}
Con el tartrato de dietilo L-($+$), (2\textit{S},3\textit{S})-2,3-epoxihexan-1-ol; con el tartrato de dietilo D-($-$), el
enantiómero, (2\textit{R},3\textit{R}).
\end{solution}

\begin{solution}{exo:b3:asymmetric-synthesis:9}
L = fenilo, M = metilo, S = H. Para el aldehído (\textit{S}), el confórmero reactivo tiene el fenilo
perpendicular al plano del carbonilo y el hidrógeno junto a la trayectoria del
nucleófilo, que ataca en anti respecto del fenilo. Al deducir las configuraciones se obtiene el
(2\textit{S},3\textit{S})-3-fenilbutan-2-ol como diastereómero mayoritario.
\end{solution}

\begin{solution}{exo:b3:asymmetric-synthesis:10}
Al principio, los dos enantiómeros están presentes en cantidades iguales, de modo que los productos
se forman en la razón $k_{\text{rápida}} : k_{\text{lenta}} = s$: $\mathrm{ee}_p = (s - 1)/(s + 1)$, 0.905 para $s = 20$. A medida que avanza la reacción, el
enantiómero rápido se agota y el ee del producto disminuye, mientras que el ee del sustrato aumenta
hacia 1: el sustrato siempre puede purificarse avanzando más; el producto
no.
\end{solution}

\begin{solution}{exo:b3:asymmetric-synthesis:11}
Con racemización rápida: hasta un 100\,\% de rendimiento, $\mathrm{er} = 99$, ee del 98\,\%. Sin ella, al 50\,\% de
conversión: un rendimiento de como máximo el 50\,\% y un ee del producto del 93\,\% (el ee del sustrato es el
mismo, 93\,\%).
\end{solution}

\begin{solution}{exo:b3:asymmetric-synthesis:12}
Cuando los dos complejos se interconvierten más deprisa de lo que reaccionan, la razón de productos es
$k_B[\mathrm B]/(k_A[\mathrm A]) = 600/10 = 60$: el complejo minoritario B da el 98\,\% del producto (ee del 97\,\%). La
selectividad la fija la diferencia entre los dos estados de transición, no
las \omterm{def:b3:partition-functions:boltzmann}{poblaciones} de los intermedios.
\end{solution}

\begin{solution}{pb:b3:asymmetric-synthesis:1}
\textbf{1.} El carbono del alqueno que lleva el nitrógeno tiene tres sustituyentes distintos
y dos caras diferentes; el hidrógeno añadido a él convierte el carbono $\alpha$ del
aminoácido en un estereocentro.

\textbf{2.} El racémico: ambas caras son atacadas a la misma velocidad.

\textbf{3.} Su oxígeno carbonílico se une al rodio junto con el enlace C=C: el sustrato queda
sujeto como un quelato, en una geometría fija que el ligando puede discriminar.

\textbf{4.} Construye un bolsillo quiral alrededor del metal (de simetría $C_2$): las dos maneras de
unir el sustrato, por una cara o por la otra, son diastereoméricas, con energías y
reactividades distintas.

\textbf{5.} No: el complejo minoritario reacciona mucho más deprisa con el hidrógeno (Curtin--Hammett).

\textbf{6.} El metal está en una cara del alqueno unido, y los hidrógenos se
transfieren desde el metal.

\textbf{7.} $\mathrm{er} = 1.95/0.05 = 39$.

\textbf{8.} $\Delta\Delta G^\ddagger = RT\ln 39 = 8.314 \times 298.15 \times 3.664 = \qty{9.1}{kJ/mol}$.

\textbf{9.} $\mathrm{er} = \eu^{9082/(8.314 \times 195.15)} = 270$: ee 99.3\,\%.

\textbf{10.} $\mathrm{er} = \eu^{11082/(8.314 \times 298.15)} = 87$: ee 97.7\,\%.

\textbf{11.} Menos que un enlace de hidrógeno típico: una pequeña diferencia en los contactos
estéricos o una sola interacción débil entre el sustrato y el ligando decide el
resultado, por lo que estos catalizadores son difíciles de diseñar a partir de primeros principios.

\textbf{12.} La velocidad disminuye, el hidrógeno se disuelve y el catalizador se comporta de otro modo
a baja temperatura, y enfriar un reactor grande cuesta energía.

\textbf{13.} 97.5 g del enantiómero (\textit{S}) y 2.5 g del (\textit{R}).

\textbf{14.} $97.5 - 5.0 = \qty{92.5}{g}$ de cristales, prácticamente enantiopuros.

\textbf{15.} $92.5/100 = 92.5$\,\% del producto bruto (el 95\,\% del enantiómero (\textit{S})).

\textbf{16.} Los cristales recogen el enantiómero en exceso, y la parte racémica queda en
disolución con el minoritario; a partir de un racémico no hay exceso que recoger, y la
cristalización da cristales racémicos.

\textbf{17.} Por HPLC quiral frente a una referencia racémica.

\textbf{18.} El 50\,\%: solo sobrevive el enantiómero lento.

\textbf{19.} $s = \ln(0.45 \times 0.05)/\ln(0.45 \times 1.95) = -3.79/-0.131 = 29$.

\textbf{20.} El 45\,\% del racémico (el 88\,\% del enantiómero deseado presente al principio,
con un ee del 95\,\%).

\textbf{21.} Con una racemización rápida del sustrato, el rendimiento podría acercarse al 100\,\%.

\textbf{22.} Convierte todo el precursor aquiral barato en el enantiómero correcto,
con un catalizador usado en cantidades ínfimas y sin desechar ningún enantiómero.

\textbf{23.} El rodio es tóxico y sus residuos en los fármacos están limitados a unas pocas partes por
millón; se miden por espectrometría de masas con plasma de acoplamiento inductivo.

\textbf{24.} Un ee del 95\,\% a \qty{25}{\celsius} corresponde a $\Delta\Delta G^\ddagger = RT\ln 39 \approx \qty{9.1}{kJ/mol}$.
\end{solution}
